% TOP_PROBLEM: {"problemId": "problem.kadison-similarity-problem-for-c-star-algebras", "canonicalTitle": "Kadison Similarity Problem for C*-Algebras", "releaseRank": 252, "primaryDomain": "algebra_representation_category", "primaryDomainLabel": "Algebra, representation & category theory", "displayStatus": "open", "releaseStatus": "open"}
% !TeX program = lualatex
% ATTEMPT_STATUS: unresolved
% Research continuation by GPT_6_Astra_Ultra, 27 September 2026.
\documentclass[11pt]{article}
\usepackage[margin=25mm]{geometry}
\usepackage{fontspec}
\setmainfont{Latin Modern Roman}
\usepackage{amsmath,amssymb,mathtools}
\usepackage{unicode-math}
\setmathfont{Latin Modern Math}
\usepackage{xurl,hyperref}
\hypersetup{colorlinks=true,urlcolor=blue}
\setcounter{secnumdepth}{0}
\setlength{\parindent}{0pt}
\setlength{\parskip}{5pt}
\setlength{\emergencystretch}{3em}
\begin{document}
\section*{\texorpdfstring{Kadison Similarity Problem for C*-Algebras}{Kadison Similarity Problem for C*-Algebras}}\label{252}
\subsection{Definitions and mathematical statement}
Let $A$ be a complex $C^*$-algebra and $H$ a Hilbert space. A bounded representation is a bounded algebra homomorphism $\pi:A\to B(H)$, not initially required to preserve adjoints. It is similar to a $*$-representation if there is a bounded invertible operator $S\in B(H)$ such that
\[
 a\longmapsto S\pi(a)S^{-1}
 \quad\text{satisfies}\quad
 S\pi(a^*)S^{-1}=(S\pi(a)S^{-1})^*.
\]
The question asks whether this holds for every bounded representation, in the unital/nondegenerate conventions of the cited formulation or the standard unitized version. The similarity operator may depend on $\pi$; no universal bound on $\|S\|\|S^{-1}\|$ is required.
\par\smallskip\textit{Formulation and concept references:} \href{https://gupea.ub.gu.se/server/api/core/bitstreams/57159657-f889-4758-943d-0b64466e9f64/content}{[S1]}.

\subsection{Short English statement}
Can every bounded representation of a C*-algebra be changed by a bounded invertible similarity so that it respects adjoints?
\subsection{Research attempt: invariant metrics and matrix amplification}
\textit{GPT-6 Astra Ultra, 27 September 2026. Status: UNRESOLVED.}

I work first with a unital homomorphism $\pi$ and write $M=\|\pi\|\ge1$. The aim is to construct a positive operator $Q$, bounded above and bounded away from zero, such that
\[
 \pi(u)^*Q\pi(u)=Q\qquad(u\in U(A)).                         \tag{1}
\]
This converts similarity into a concrete averaging problem. It gives a full proof for finite-dimensional algebras and for a specified class of norm limits. It also exposes why boundedness at the scalar level does not by itself provide the matrix estimates required in general.

\paragraph{From an invariant metric to a star representation.}
Assume $cI\le Q\le CI$ for $0<c\le C<\infty$ and (1). Let $S=Q^{1/2}$. Then $S\pi(u)S^{-1}$ is an invertible isometry, hence unitary, for every unitary $u\in A$. A self-adjoint contraction $a$ is the real part of a unitary:
\[
 u=a+i(1-a^2)^{1/2},\qquad a=(u+u^*)/2.
\]
Since $u^*=u^{-1}$, the homomorphism $S\pi(\cdot)S^{-1}$ carries $u^*$ to the adjoint of the image of $u$. It therefore carries $a$ to a self-adjoint operator. Decomposing an arbitrary element into its real and imaginary self-adjoint parts proves preservation of adjoints. Thus (1), with the two-sided bound, suffices. Conversely, a similarity $S$ to a star representation gives such a metric $Q=S^*S$. No compactness or finite dimensionality is involved in this equivalence.

The lower bound on $Q$ is essential. An invariant positive semidefinite form with a kernel would only define a representation on a quotient. A positive injective operator whose spectrum accumulates at zero would give an unbounded inverse square root, which is also outside the problem's conclusion.

\paragraph{A complete averaging proof for finite-dimensional $A$.}
For a finite-dimensional $C^*$-algebra, its unitary group is compact. With normalized Haar measure set
\[
 Q=\int_{U(A)}\pi(u)^*\pi(u)\,du.
\]
Because $\pi(u)^{-1}=\pi(u^*)$ and both have norm at most $M$,
\[
 M^{-2}\|\xi\|^2\le\|\pi(u)\xi\|^2\le M^2\|\xi\|^2.
\]
Integration gives $M^{-2}I\le Q\le M^2I$. Right invariance of Haar measure gives, for $v\in U(A)$,
\[
 \pi(v)^*Q\pi(v)
 =\int\pi(uv)^*\pi(uv)\,du=Q.
\]
Consequently $S=Q^{1/2}$ works and
\[
 \|S\|\,\|S^{-1}\|\le M^2.                                \tag{2}
\]
The representation space $H$ can be infinite-dimensional: only the compactness of $U(A)$ was used. The integral can be read weakly through matrix coefficients, so no assumption that $B(H)$ is finite-dimensional is hidden here.

This proof also handles a nonunital bounded homomorphism from a unital algebra after an initial similarity. The idempotent $P=\pi(1)$ gives the topological direct sum $H=\operatorname{ran}P\oplus\ker P$. An equivalent Hilbert norm makes this decomposition orthogonal. The representation is unital on its range and zero on its kernel. This reduction does not say that $P$ was originally an orthogonal projection; requiring that would already assume part of the desired conclusion.

\paragraph{A limit argument with a genuine uniform bound.}
Suppose $A$ is unital and contains a directed family of finite-dimensional unital star subalgebras $A_i$ whose union is norm dense. Apply the preceding construction to $\pi|_{A_i}$. The resulting $Q_i$ all lie in the same weak-operator compact order interval
\[
 M^{-2}I\le Q_i\le M^2I.
\]
Choose a convergent subnet with limit $Q$. For a fixed $u\in U(A_j)$, the relation $\pi(u)^*Q_i\pi(u)=Q_i$ holds whenever $i\ge j$; multiplication on either side by a fixed bounded operator is weak-operator continuous. Thus $Q$ satisfies the same relation for every unitary in every $A_j$.

It follows that the conjugated homomorphism preserves adjoints on each $A_j$. Norm continuity and density extend this identity to all of $A$. The bounds on $Q$ survive the limit, so $Q^{1/2}$ is boundedly invertible and (2) still holds. This proves the conclusion for algebras having precisely this approximation property, including unital AF algebras.

There is a useful lesson in the order of this argument. It is not enough that each finite collection of equations admits some positive solution. The constants must stay away from zero and infinity in order to take a limit. Nor is a general $C^*$-algebra known to have the finite-dimensional subalgebra approximation used above. Completely positive approximations, when available, are a different object; their images are not automatically subalgebras on which the same averaging argument can be performed.

\paragraph{The matrix obstruction can be stated sharply.}
For $n\ge1$ let $\pi_n:M_n(A)\to M_n(B(H))$ act entrywise. If $\pi$ is similar to a star representation through $S$, then amplification conjugates by $I_n\otimes S$, so
\[
 \|\pi_n\|\le\|S\|\,\|S^{-1}\|\quad\hbox{for every }n.
                                                               \tag{3}
\]
In particular $\|\pi\|_{\mathrm{cb}}:=\sup_n\|\pi_n\|$ is finite. Haagerup's similarity theorem gives the converse for bounded homomorphisms of $C^*$-algebras [E1]. Using that theorem, the remaining problem is exactly whether every bounded algebra homomorphism has finite completely bounded norm. The present proof establishes this finiteness for the approximation class above; it does not establish it for arbitrary $A$.

A direct estimate illustrates the missing uniformity. If $X=[a_{ij}]\in M_n(A)$, then $\|a_{ij}\|\le\|X\|$. For $\xi=(\xi_j)$, Cauchy--Schwarz gives
\[
 \|\pi_n(X)\xi\|^2
 \le nM^2\|X\|^2\Big(\sum_j\|\xi_j\|\Big)^2
 \le n^2M^2\|X\|^2\|\xi\|^2.
\]
Thus $\|\pi_n\|\le nM$. Each individual amplification is bounded, but taking the supremum of this estimate gives nothing. There is no logical implication from ``finite for every $n$'' to ``bounded uniformly in $n$.'' Multiplicativity must enter a stronger argument somewhere.

\paragraph{An adversarial test for scalar-norm arguments.}
Let $T:M_d(\mathbb C)\to M_d(\mathbb C)$ be transpose. It is a linear isometry. On $\mathbb C^d\otimes\mathbb C^d$, the flip
\[
 F=\sum_{i,j}E_{ij}\otimes E_{ji}
\]
is unitary, whereas
\[
 (\mathrm{id}\otimes T)(F)
 =\sum_{i,j}E_{ij}\otimes E_{ij}
 =|\Omega\rangle\langle\Omega|,\qquad
 \Omega=\sum_i e_i\otimes e_i,
\]
has norm $d$. Therefore scalar norm one does not control amplification even in a finite-dimensional matrix space. This is explicitly \emph{not} a counterexample to Kadison: transpose reverses products and is not an algebra homomorphism when $d>1$. Its role is to reject any attempted proof that uses only the norm estimate and never uses the order of multiplication.

There is a parallel warning about averaging. Compact Haar averaging on $U(A_i)$ is legitimate in the finite-dimensional proof. For general $A$, its unitary group is usually not compact in the norm topology. Declaring a normalized invariant average on all its bounded functions is an additional amenability assumption, not a consequence of the existence of the norm bound $M$. One cannot insert that average as a formal replacement for the integral without proving its existence in the required setting.

\paragraph{Verification and remaining gap.}
The local exact checks verify the flip and rank-one identities for several $d$, and verify the metric equation on an explicit nonorthogonal two-dimensional representation. They test the algebra and the direction of the similarity, not the missing infinite-dimensional theorem. The established result is constructive similarity for finite-dimensional algebras and for the directed approximation class above. The first unresolved step is obtaining a uniformly coercive invariant metric, or equivalently the finiteness in (3), for every bounded homomorphism of an arbitrary $C^*$-algebra. Neither the compactness argument nor the scalar estimates bridge that step here.

\subsection{Sources}
\begin{itemize}
\item[S1]Tim Westlund, Similarity Problems: Which Groups Are Unitarizable?, Master's thesis, University of Gothenburg, 2024.
\url{https://gupea.ub.gu.se/server/api/core/bitstreams/57159657-f889-4758-943d-0b64466e9f64/content}.
\textit{Formulation source.}
\item[E1]U. Haagerup, \textit{Solution of the similarity problem for cyclic representations of $C^*$-algebras}, Annals of Mathematics 118 (1983), 215--240. The completely bounded similarity criterion and cyclic-representation theorem are external results; the averaging and limit arguments above are written out separately.
\url{https://annals.math.princeton.edu/1983/118-2/p01}.

\end{itemize}
\end{document}
